% Author-review candidate research subsection only.
% Not included in the canonical journal manuscript, because replication
% beyond one public acquisition family and known synchronisation baselines
% has not established a separate original algorithm contribution.
\subsection{Controlled real-signal test of coherent multi-sensor release provenance}
\label{sec:p2-source-clock-positive}
We distinguish \emph{any individually reused channel samples} from
\emph{one coherently time-synchronised multichannel recording segment}.
The latter is the explicitly narrower target of this positive mechanism
experiment. It requires at least four distinct sensor correspondences
to support a single source--target time offset. Candidate correlations
are screened at absolute Pearson correlation $\geq0.985$ with at least
230 overlapping samples and permitted offsets within $\pm100$ samples.
One-to-one channel correspondences are resolved by maximum-weight
assignment. These settings were held fixed before examining the new
S005--S006 test IDs.

Public PhysioNet EEGMMIDB EDF v1.0.0 signals were used as the source
waveforms. Only the first eight EEG sensors and 768 samples were used
per ID/record. IDs S001--S002 were development sources. S003--S004
were tested previously as a held proof-of-concept; a later revision
therefore used previously uninspected S005--S006 as the additional held
ID pair. Every source yielded eight \emph{experimentally controlled}
same-waveform positive variants (including channel renaming, cropping,
fixed lag, scaling, channel omission, and low-amplitude noise) and
four constructed negative conditions. This is not naturally observed
clinical duplication ground truth.

\begin{table}[tb]
\centering
\small
\caption{Observed results on the new S005--S006 controlled-transform
held set (two public source IDs). All numbers are case counts, not
independent clinical subject counts. The last negative deliberately
reuses all source channels at different delays and is negative
\emph{only for common-clock multichannel ancestry}.}
\label{tab:p2-source-clock-physionet-new-held}
\begin{tabular}{lrr}
\toprule
Method & Detected positives / 16 & False matches / 8\\
\midrule
Shared-clock, one-to-one channels & 16 & 0\\
Independent-best-lag channel Hungarian & 16 & 2\\
Zero-lag channel Hungarian & 10 & 0\\
Named-channel shifted Pearson & 12 & 0\\
Named-channel zero-lag Pearson & 8 & 0\\
Full labelled numerical equality & 2 & 0\\
\bottomrule
\end{tabular}
\end{table}

For the independent-best-lag Hungarian baseline, the two false matches
were the constructed \texttt{asynchronous\_same\_source} scenarios:
each matched sensor was digitally copied from the same source, but
distinct per-channel delays prevented any four sensor traces from
sharing a single clock. The coherent matcher rejected both without
losing any of the 16 constructed held positives. This controlled
contrast corroborates the \emph{operational utility of a shared-time
constraint}, not novelty over established EEG synchronization methods,
nor a general improvement in any-channel digital content reuse
detection. The earlier S003--S004 controlled test showed the same
16/16 positives and 0/8 versus 2/8 false-match pattern. The source
EDFs and all post-transform outcomes remain publicly traceable in
independent GitHub Actions receipts
\url{https://github.com/lindicaphxag-tech/kaggle/actions/runs/38662768329}
and
\url{https://github.com/lindicaphxag-tech/kaggle/actions/runs/38663141073}.

The experiment uses one open EEG acquisition family, very few distinct
held people, intentionally constructed negatives, and a contrast
whose preferred label depends on the intended meaning of provenance.
Consequently it does \emph{not} establish sensitivity or specificity
for clinical reuse, clinical participant identity, original MAT--EDF
ancestry, or domain-independent release auditing. Additional independent
EEG datasets and comparisons with contemporary multichannel
synchronization, approximate fingerprinting and time-warping methods
are required before an original-algorithm superiority claim.
